\documentclass[12pt]{article}
\usepackage{polyglossia}
\usepackage{fontspec}
\usepackage[export]{adjustbox}
\usepackage{amsfonts, amsmath, amssymb}
\usepackage{annotate-equations}
\usepackage{bm}
\usepackage{bookmark}
% \usepackage{breqn}
\usepackage{chemformula}
\usepackage{derivative}
\usepackage{empheq}
\usepackage[inline]{enumitem}
\usepackage{fancyhdr}
\usepackage{geometry}
\usepackage{graphicx}
\usepackage{kantlipsum}
\usepackage{lastpage}
\usepackage{linebreaker} % line-breaker algorithm in LuaLaTeX
\usepackage{lualatex-math} % Fixes for mathematics
\usepackage{mathtools}
\usepackage{microtype}
\usepackage{mismath}
\usepackage{nicefrac}
\usepackage[hyperref]{ntheorem}
% \usepackage{phfqit} % Bra-ket notation
% \usepackage{physics}
\usepackage[final]{showlabels}
\usepackage{simples-matrices}
\usepackage{siunitx}
\usepackage{soul}
\usepackage{spbmark}
\usepackage{titling}
\usepackage{tabularray}
\usepackage{tasks}
\usepackage{tensor}
\usepackage{xcolor}
\usepackage{xfrac}
\usepackage{xparse}
\usepackage{xspace}
\usepackage{hyperref}
% Polyglossia config
\setdefaultlanguage[variant=mexican]{spanish}
\setotherlanguage{english}
\usepackage{cleveref}

%%%%%%%%%%%%%%%%%%%%%%%%%%%%%%%%%%%%%%%%%%%%%%%%%%%%%%%%%%
%%%%%%%%%%%%%%%%%%%%%%%%%%%%%%%%%%%%%%%%%%%%%%%%%%%%%%%%%%
%%%%%%%%%%%%%%%%%%Configuraciones extras%%%%%%%%%%%%%%%%%%
\definecolor{base3}{RGB}{253, 246, 227}%
\definecolor{pinkwave}{RGB}{255, 0, 128}%
\definecolor{customBlue}{RGB}{11, 61, 98}%
\pagecolor{base3}
\graphicspath{{img/}}
\setlength{\parindent}{2em} % Sangría
\setlength{\parskip}{0.5em} % Espacio entre párrafos
\linespread{1.1} % line spacing
\setlength{\jot}{10pt} % Space between lines in multiline eqs
% \crefname{equation}{ec.}{ecs.} % Equation's cross-reference name
\crefname{equation}{}{} % Equation's cross-reference name

% Line-breaker config
\linebreakersetup {
    maxtolerance = 90,
    maxemergencystretch = 1em,
    maxcycles = 4
}

%%%%%%%%%%%%%%%%%%Theorem environments%%%%%%%%%%%%%%%%%%
% Configuración de ambiente para problema
\theoremstyle{break}
\theoremheaderfont{\Large\normalfont\bfseries}
\theorembodyfont{\normalfont}
\theoremseparator{\bigskip} % Spacing between header and body
\theorempreskip{1.5em}
\theorempostskip{\topsep\bigskip}
\theorempostwork{
    \color{customBlue} \hrule width \hsize height 2pt \kern 1mm \hrule width \hsize
    }
\newtheorem{exercise}{Problema}
\counterwithin{equation}{exercise}

\makeatletter
\LetLtxMacro\old@exercise\exercise
\renewcommand*{\exercise}{%
  \@ifnextchar[{%
    \exercise@opt
  }{%
    \old@exercise
    \exercise@bookmark{}%
  }%
}
\def\exercise@opt[#1]{%
  \old@exercise[{#1}]%
  \exercise@bookmark{: #1}%
}
\newcommand*{\exercise@bookmark}[1]{%
  \bookmark[dest=\@currentHref]{Problema \theexercise#1}%
}
\makeatother

% Configuración de ambiente para solución
\theoremstyle{nonumberbreak}
\theoremheaderfont{\Large\normalfont\bfseries}
\theorembodyfont{\normalfont}
\theoremseparator{\medskip}
\theorempreskip{1em}
\theorempostskip{\topsep\medskip}
\newtheorem{solution}{Solución}

%%%%%%%%%%%%%%%%%%%%%%%%%%%%%%%%%%%%%%%%%%%%%%%%%%%%%%%%


% Configuración del paquete hyperref
\hypersetup{
    colorlinks = true,%
    linkcolor={[rgb]{0,0.2,0.6}},%
    citecolor={[rgb]{0,0.6,0.2}},%
    filecolor={[rgb]{0.8,0,0.8}},%
    urlcolor={[rgb]{0.8,0,0.8}},%
    runcolor={[rgb]{0.8,0,0.8}},% 
    menucolor={[rgb]{0,0.2,0.6}},%
    linkbordercolor={[rgb]{0,0.2,0.6}},%
    citebordercolor={[rgb]{0,0.6,0.2}},%
    filebordercolor={[rgb]{0.8,0,0.8}},%
    urlbordercolor={[rgb]{0.8,0,0.8}},%
    runbordercolor={[rgb]{0.8,0,0.8}},%
    menubordercolor={[rgb]{0,0.2,0.6}},% 
    pdftitle={Tarea 2},%
    pdfauthor={López Merino Marcos},%
    pdfsubject={Introducción a la Física Nuclear},%
    pdfkeywords={Facultad de Ciencias, UNAM, Introducción a la Física Nuclear, Experimental},%
    unicode = true%
}

%%%%%%%%%%%%%%%%%%siunitx configuration%%%%%%%%%%%%%%%%%%
% Configuración del paquete siunitx
\sisetup{
	output-decimal-marker = {.}, 
	per-mode = symbol-or-fraction,
	separate-uncertainty = false,
	exponent-product = \mul,
    % inter-unit-product = \ensuremath{{}\vdot{}}
}

% Declaring new units
\DeclareSIUnit\kilogram{\kilo\gram}
\DeclareSIUnit\nucleones{\text{nucleones}}
\DeclareSIUnit\density{\gram\per\cm\cubed}

%%%%%%%%%%%%%%%%%%%%%%%%%%%%%%%%%%%%%%%%%%%%%%%%%%%%%%%%

% Geometría del documento
\geometry{
    letterpaper,
    top = 0.6in,
    bottom = 0.8in,
    left = 0.6in,
    right = 0.6in,
    footskip = 38pt
}

%%%%%%%%%%%%%%%%%%Nuevos comandos%%%%%%%%%%%%%%%%%%
\newcommand*{\group}{8316}
\newcommand*{\classname}{Introducción a la Física Nuclear}
\newcommand*{\homeworknumber}{Tarea 2}
\newcommand*{\name}{Marcos López Merino}

% unit vector i
\newcommand*{\uveci}{{\bm{\hat{\textnormal{\bfseries\i}}}}}
% unit vector j
\newcommand*{\uvecj}{{\bm{\hat{\textnormal{\bfseries\j}}}}}
% unit vector
\DeclareRobustCommand{\uvec}[1]{{%
  \ifcsname uvec#1\endcsname
     \csname uvec#1\endcsname
   \else
    \bm{\hat{\mathbf{#1}}}%
   \fi
}}% 
\newcommand{\idest}{\emph{i.e.},\xspace} % id est
% Espacio vectorial, e.g., ℝ, ℂ, ℕ, etc.
\NewDocumentCommand{\vecspace}{m o}{%
  \IfValueTF{#2}{%
    \mathbb{#1}^{#2}%
  }{%
    \mathbb{#1}%
  }%
}

% unit vector k
\newcommand*{\uvectk}{\uvec{k}}

\newcommand*{\observer}{\mathcal{O}}
\newcommand*{\primeobserver}{\mathcal{O}^{\prime}}
\newcommand*{\biprimeobserver}{\mathcal{O}^{\prime\prime}}
\newcommand*{\triprimeobserver}{\mathcal{O}^{\prime\prime\prime}}
\newcommand*{\inlinesol}{\vspace*{10pt}\textbf{Solución}\vspace*{10pt}}
\renewcommand*{\grad}{\nabla}
\renewcommand*{\div}{\nabla \cdot}
\renewcommand*{\curl}{\nabla \mul}
\newcommand*{\laplacian}{\nabla^{2}}
\newcommand*{\qq}{\quad\quad}
\DeclareDifferential{\cdif}{\grad}[style-notation=single]
\newcommand*{\dotprod}[2]{#1 \mathbin{\cdot} #2\xspace}
\newcommand*{\crossprod}[2]{#1 \mathbin{\mul} #2\xspace}
\newcommand*{\basevect}[1]{\uvec{e}_{#1}\xspace}

% \newcommand*{\change}[1]{\Delta{#1}\xspace}

\NewDocumentCommand{\change}{ s o m }{%
    \IfBooleanTF{#1}{%
        \Delta {#3^{\prime}}
    }{%
        \Delta {#3}
    }
}
\NewDocumentCommand{\interval}{ s o m }{%
    \IfBooleanTF{#1}{%
        (\Delta {#3^{\prime}})^{2}
    }{%
        (\Delta {#3})^{2}
    }
}
\DeclarePairedDelimiterX\set[1]\lbrace\rbrace{#1}
% \newcommand{\crefrangeconjunction}{--}
% \newcommand{\crefpairconjunction}{\xspace y\xspace}

% Defining a variant of Aboxed command from mathtools
\makeatletter
\newcommand*\Acolorboxed[2][pinkwave]{%
   \let\bgroup{\romannumeral-`}%
   \@Acolorboxed{#1}#2&&\ENDDNE
}
\def\@Acolorboxed#1#2&#3&#4\ENDDNE{%
  \ifnum0=`{}\fi
  \setbox\z@\hbox{$\displaystyle#2{}\m@th$\kern\fboxsep \kern\fboxrule}%
  \edef\@tempa{\kern\wd\z@ & \kern-\the\wd\z@ \fboxsep\the\fboxsep \fboxrule\the\fboxrule}%
  \@tempa
  \fcolorbox{#1}{base3}{\m@th$\displaystyle#2#3$}%
} 
\makeatother

%%%%%%%%%%%%%%%%%%Portada y configuración%%%%%%%%%%%%%%%%%%
% Configuración de portada
\setlength{\droptitle}{-60pt} % raise the title

% Portada
\title{
    \textbf{\homeworknumber}\\
    \normalsize\vspace{0.1in}\small{\textbf{Entrega}:~\today}
    \vspace{-1.5in}
}
\author{}
\date{}

%%%%%%%%%%%%%%%%%%Header and footer%%%%%%%%%%%%%%%%%%
\setlength{\headheight}{15.2pt}
\pagestyle{fancy}
\lhead{Grupo \group, Sem. 2023-1}
\chead{\classname}
\rhead{\name}
\lfoot{\includegraphics[scale = 0.2, valign = c]{LogoFCUNAMcolor.pdf}}
% \lfoot{\includegraphics[scale = 0.1, valign = c]{example-image}}
\cfoot{\homeworknumber}
\rfoot{Pág. \thepage \hspace{1pt} de \pageref{LastPage}}

\renewcommand{\headrulewidth}{0.5pt}
\renewcommand{\footrulewidth}{0.5pt}
%%%%%%%%%%%%%%%%%%%%%%%%%%%%%%%%%%%%%%%%%%%%%%%%%%%%%%%%%%
%%%%%%%%%%%%%%%%%%%%%%%%%%%%%%%%%%%%%%%%%%%%%%%%%%%%%%%%%%

\begin{document}
    \maketitle
    \thispagestyle{fancy}

    % Problema 1
    \begin{exercise}
      Calcular la masas total del planeta Tierra con los siguientes datos:

      Suponiendo que se trata de una esfera de radio \(r = \SI{6400}{\km}\); y que además la Tierra está constituida por los siguientes porcentajes de elementos: \SI{37}{\percent} de \ch{Fe}, \SI{25}{\percent} de \ch{Si}, \SI{13}{\percent} de \ch{Mg}, \SI{10}{\percent} de \ch{Ni}, \SI{8}{\percent} de \ch{Ca} y \SI{7}{\percent} de \ch{K}.

      Se deben consultarlas densidades de los componentes: \(\rho = \left[\unit[per-mode=fraction]{\density}\right]\) que se requieren para calcular la masa total de la Tierra.

      \begin{solution}
        Las densidades de los elementos se obtuvieron de \href{https://www.physics.nist.gov/cgi-bin/Star/compos.pl}{Material composition data}:

        \begin{alignat*}{3}
          \rho_{\ch{Fe}} &= \SI{7874}{\density},\quad \rho_{\ch{Si}} &{}={}& \SI{2330}{\density},\quad \rho_{\ch{Mg}} &{}={}& \SI{1740}{\density},\\
        \rho_{\ch{Ni}} &= \SI{8902}{\density},\quad \rho_{\ch{Ca}} &{}={}& \SI{1550}{\density},\qquad \rho_{\ch{K}} &{}={}& \SI{862}{\density}.
        \end{alignat*}

        Y como queremos calcular la masa total de la Tierra, entonces la densidad total es:

        \begin{equation}
          \rho = 0.37 \rho_{\ch{Fe}} + 0.25 \rho_{\ch{Si}} + 0.13 \rho_{\ch{Mg}} + 0.10 \rho_{\ch{Ni}} + 0.08 \rho_{\ch{Ca}} + 0.07 \rho_{\ch{K}}.
          \label{eq:TotalDensity}
        \end{equation}

        Sustituyendo los valores de las densidades en \cref{eq:TotalDensity}

        \begin{align}
          \rho &= 0.37\left(\SI{7874}{\density}\right) + 0.25\left(\SI{2330}{\density}\right) + 0.13\left(\SI{1740}{\density}\right),\nonumber\\
          &+ 0.10\left(\SI{8902}{\density}\right) + 0.08\left(\SI{1550}{\density}\right) + 0.07\left(\SI{862}{\density}\right),\nonumber\\
          \Acolorboxed{\rho &= \SI{4796.62}{\density}.}\label{eq:TotalDensityN}
        \end{align}

        Y, además, la masa total de la Tierra está dada por:

        \begin{align}
          M &= \rho V,\nonumber\\
          M &= \rho \left(\dfrac{4}{3}\pi r^{3}\right),\label{eq:Mass}
        \end{align}

        donde \(\frac{4}{3}\pi r^{3}\) es el volumen de una esfera.

        Sustituyendo \cref{eq:TotalDensityN} y el radio de la Tierra \(r = \SI{6.4e8}{\cm\cubed}\) en \cref{eq:Mass}:

        \begin{align*}
          M &= \SI{4798.68}{\density} \left(\dfrac{4}{3}\pi \left(\SI{6.4e8}{\cm}\right)^{3}\right),\\
          &= \SI{5.34886e30}{\gram},\\
          \Acolorboxed{M &= \SI{5.34886e27}{\kg}.}
        \end{align*}
      \end{solution}
    \end{exercise}

    % Problema 2
    \begin{exercise}
      Por otro lado, se puede ilustrar que la materia a nivel microscópico es hueva. Hacer el siguiente cálculo:

      Suponer que se tiene un balín esférico de radio \(r = \SI{1}{\cm}\) compuesto de \ch{^{nat} Fe} (con \(A = 54\) (\SI{6}{\percent}), \(A = 56\) (\SI{92}{\percent}) y \(A = 57\) (\SI{2}{\percent}), isótopos más abundantes del \ch{Fe}). Calcular el volumen de un núcleo de \ch{Fe} cuyo radio es \(r = r_{0}A^{\sfrac{1}{3}}\). Suponiendo que no hay repulsión coulombiana, ¿cuántos átomos de \ch{Fe} cabrían en el balín de \SI{1}{\cm} de radio? Calcular en [\unit{\kg}] lo que pesaría el balín con esa cantidad de átomos.

      \textbf{NOTA}: Tomen el valor de \(r_{0}\) con las unidades convenientes.
    \end{exercise}

    % Problema 3
    \begin{exercise}
      De tus resultados anteriores, ¿cuántos balines son necesarios para compararlos con el peso de la Tierra? ¿Qué se puede concluir al respecto?
    \end{exercise}
\end{document}